\documentclass{ximera}

\usepackage{pgfplots}
\pgfplotsset{compat=1.18}


\newcommand{\rulecolor}[1]{\def\CT@arc@{\color{#1}}}
\definecolor{darkblue}{rgb}{0.4,0.5,0.85}
\definecolor{blue}{rgb}{0.1,0.3,.95}
%\definecolor{blue}{rgb}{0.1,0.2,0.90}
\definecolor{pink}{rgb}{0.6,0.4,0.7}
\definecolor{green}{rgb}{0.4,0.6,0.5}
\definecolor{yellow}{rgb}{0.1,0.6,0.8}
\definecolor{red}{rgb}{0.8,0.1,0.5}
\definecolor{glaucous}{rgb}{0.38,0.51,0.71}
\definecolor{darkraspberry}{rgb}{0.53, 0.15, 0.34}


\title{A Basic Review of Integration}   % Use xmlatex all -s to compile when SERVE refuses to work.
\author{Kelly Stady}
\license{CC: 0}         % replace with an appropriate license, or set it in xmPreamble

\begin{document}
\begin{abstract}
    A basic review of antiderivatives for students who have practiced.......
\end{abstract}
\maketitle
\label{xim:IntegrationActivity}



\section*{Antiderivatives}

We will now reverse the process of differentiation.  That is, we will be
given a derivative and asked to find the original function. The original function is called
the \textbf{antiderivative} and the process of finding it is called
\textbf{antidifferentiation} or \textbf{integration}.

\begin{theorem}{Definition of an Antiderivative}
A function $F(x)$ is an \textbf{antiderivative} of $f(x)$ if $F'(x)=f(x).$
\end{theorem}

\begin{remark}
Not all functions have antiderivatives, but those that do have infinitely many.  For example, $F(x)=x^2$ is
one antiderivative of $f(x)=2x$ since $F'(x) = f(x).$  However,
$F(x)= x^2+C,$ where $C$ is a constant, is also an antiderivative
of $f(x)=2x,$ thus there are infinitely many antiderivatives of
$f(x)=2x.$

\end{remark}



\begin{problem}
Is $F(x) = 3x^2+5$ an antiderivative of $f(x) = 6x?$

\begin{multipleChoice}
\choice{Yes, $F'(x) = f(x).$}
\choice{No, $F'(x) \ne f(x).$}
\end{multipleChoice}


\end{problem}


\begin{problem}
Identify all antiderivatives of $f(x)=x^3+9.$

\begin{selectAll}
\choice[correct]{$F(x) = \frac{x^4}{4} + 9x -5$}
\choice{$F(x) = x^4 + 9x + 2$}
\choice[correct]{$F(x) = \frac{x^4}{4} + 9x$}
\choice{$F(x) = x^4 + 9x + 6$}
\end{selectAll}

\end{problem}


\subsection*{Notation}

The following notation is used to denote the process of integration and is read as
``the integral of $f(x)$."  The elongated s symbol $\textstyle \int$ is called the \textbf{integral sign} and
the function $f(x)$ is called the \textbf{integrand}.  The differential symbol $dx$
indicates that integration is with respect to the variable $x.$  Together, these symbols are referred to as
an \textbf{indefinite integral}.

\[ \int f(x) \ dx  \hspace{0.4in} \textrm{\textbf{Indefinite Integral}} \]


If $F(x)$ is an antiderivative of $f(x)$ then we can write

\[ \int f(x) \ dx = F(x) + C \]

where $C$ is the \textbf{constant of integration}.

\begin{problem}
Identify the integrand in the problem below.

\[ \int 5x + 4 \ dx = \frac{x^5}{5} + 4x + C \]

\begin{multipleChoice}
\choice{$dx$}
\choice[correct]{$5x+4$}
\choice{$\frac{x^5}{5} + 4x + C$}
\choice{$\textstyle \int$}
\end{multipleChoice}

\end{problem}


\begin{theorem}[\textbf{Basic Integration Rules}]

You should know the following basic integration rules. The integration formulas for $\tan x$ and $\cot x$ aren’t
listed since they are easily derived by rewriting them in terms of sine and cosine and using a $u$-substitution.

\begin{enumerate}
\item $\displaystyle \int 1 \ dx = x+C $


\item $\displaystyle \int x^n \ dx = \frac{x^{n+1}}{n+1} + C, \hspace{0.3in} n \ne -1 $


\item $\displaystyle \int \frac{1}{x} \ dx = \ln |x| + C $


\item $\displaystyle \int e^x \ dx = e^x + C $


\item $\displaystyle \int \sin x \ dx = - \cos x + C $


\item $\displaystyle \int \cos x \ dx = \sin x + C$


\item $\displaystyle  \int \sec^2 x \ dx = \tan x + C $


\item $\displaystyle \int \csc^2 x \ dx = -\cot x + C $


\item $\displaystyle \int \sec x \tan x \ dx = \sec x + C $


\item $\displaystyle \int \csc x \ \cot x \ dx = -\csc x + C $


\item $\displaystyle \int \frac{1}{\sqrt{a^2-u^2}} \ du =  \arcsin \frac{u}{a} + C, \hspace{0.2in} a > 0 $


\item $\displaystyle \int \frac{1}{a^2 + u^2} \ du =  \frac{1}{a} \arctan \frac{u}{a} + C, \hspace{0.2in} a > 0$

\end{enumerate}
\end{theorem}


\begin{remark}
Often, we must first alter the form of an integrand prior to applying a basic integration rule. If an integrand is not in a
familiar form, try using one of the following techniques to create a familiar form.


\begin{enumerate}
\item Expand a power.

\item Separate a numerator into two or more terms.

\item Use a trigonometric identity.

\item Use a $u$-substitution. Try letting $u$ equal the inner most part of a composite function.

\item Use polynomial division when the degree of the numerator is greater than or equal
to the degree of the denominator.

\item Complete the square.

\item Look for integrands of the form $\textstyle \frac{u'}{u}$ or something close to that. Then let $u$ equal the
denominator. These integrals yield the natural logarithmic function.

\item If an integrand contains a logarithmic function of the form $f(x) = \ln g(x),$ try letting $u = \ln g(x).$ It won’t work in every case, but it often will.

\item Look for integrands with denominators of the form $a^2+u^2$ and $\sqrt{a^2+ u^2}.$ These often result in an answer involving the inverse tangent function or inverse sine function.

\end{enumerate}
\end{remark}


\end{document}